% core/theory/derivation.tex
\documentclass[a4paper,11pt]{article}
\usepackage{amsmath,amssymb,bm,geometry}
\geometry{margin=2cm}
\title{Theoretical Derivation: ReMetaBayes}
\author{Aulia Octaviani}
\date{}

\begin{document}
\maketitle

\section*{1. Problem Setting}

We consider a reinforcement learning agent interacting with a non-stationary environment.
At each time step $t$, the environment dynamics are parameterized by $\theta_t$, which evolves as:
\[
p(\theta_t \mid \theta_{t-1}) = \mathcal{N}(\theta_{t-1}, \sigma^2 I)
\]
The agent maintains a probabilistic belief $q(\theta_t)$ and updates it upon observing transitions $(s_t, a_t, r_t, s_{t+1})$.

\section*{2. Bayesian Meta-Update}

We define a Bayesian meta-update rule:
\[
q(\theta_t) \propto q(\theta_{t-1}) 
\exp\Big[-\alpha \, \mathbb{E}_{(s,a,r)} [L(\pi_\phi(a|s), r)]\Big]
\]
where $L(\pi_\phi, r)$ is the expected loss under the policy $\pi_\phi(a|s)$.

This formulation links the inner-loop gradient step of meta-learning to a probabilistic inference update.

\section*{3. Adaptation Speed Bound}

We define the adaptation speed $A$ as the inverse of the time required for $q(\theta_t)$ to converge after a shift.

Using information geometry arguments, we obtain:
\[
A \le \frac{1}{\lambda} 
\sqrt{\frac{2}{KL(q(\theta_t)\,||\,q(\theta_{t-1})) + \epsilon}}
\]
This bound connects learning rate $\lambda$ and information divergence between successive beliefs.

\section*{4. Regret under Distribution Shifts}

Let $\mathcal{R}_t$ be the instantaneous regret under environment $\theta_t$:
\[
\mathcal{R}_t = V^*(\theta_t) - V^{\pi_\phi}(\theta_t)
\]
For a sequence of $T$ shifts, the cumulative regret satisfies:
\[
\mathbb{E}\Big[\sum_{t=1}^T \mathcal{R}_t\Big] 
\le C \sum_{t=1}^T \sqrt{KL(q(\theta_t)||q(\theta_{t-1}))}
\]
where $C$ is a constant related to task similarity.

\section*{5. Discussion}

The result indicates that minimizing the divergence between consecutive posteriors 
reduces cumulative regret under non-stationary dynamics. 
This provides a theoretical grounding for ReMetaBayes as a probabilistic 
framework for autonomous adaptation.

\end{document}
